\documentclass[../../../include/open-logic-section]{subfiles}

\begin{document}

\olfileid{sth}{cardinals}{cp}
\olsection{कँटरचे तत्त्व}

मागे \olref[ordinals][vn]{sec} कडे मन वळवा. आपण सुक्रमित संचांविषयी चर्चा करत होतो आणि सुक्रमणांचे प्रतिनिधित्व करणाऱ्या वस्तू असाव्यात, असे सुचवले होते. या उद्देशाने आपण क्रमसूचक संख्या मांडल्या; नंतर \olref[ordinals][ordtype]{ordtypesworklikeyouwant} मध्ये दाखवले की त्यांचे गुणधर्म अपेक्षेप्रमाणे आहेत, म्हणजे:
\[
  \ordtype{A, <} = \ordtype{B, \lessdot}
  \text{ iff } \tuple{A, <} \isomorphic \tuple{B, \lessdot}.
\]
आता आणखी मागे, \olref[sfr][siz][equ]{sec} कडे वळा. तेथे संचाच्या “आकारमाना”ची कल्पना आपण सहजरीत्या मांडली होती. विशेषतः, दोन संच तुल्यबल आहेत, म्हणजे $\cardeq{A}{B}$, असे आपण तेव्हा आणि तेव्हाच म्हणालो जेव्हा $f \colon A \to
B$ हे !!a{bijection} असते. सुक्रमणाच्या कल्पनेपेक्षा ही स्वभावतः सोपी कल्पना आहे: आपल्याला फक्त !!{bijection}s महत्त्वाची आहेत; क्रम-समरूपणांच्या बाबतीत जसे पाहतो तसे, ती !!{bijection}s कोणतीही रचना जतन करतात का, हे येथे पाहायचे नाही.

यावरून एक स्वाभाविक विचार सुचतो. सुक्रमणांचे मोजमाप करण्यासाठी जशा काही वस्तू, म्हणजे \emph{क्रमसूचक संख्या}, आपण मांडल्या, तशाच आकारमानाचे मोजमाप करण्यासाठी \emph{संचांक} मांडू शकतो. हेच या प्रकरणाचे उद्दिष्ट आहे.

हे संचांक नेमके काय असतील ते सांगण्याआधी, त्यांनी पूर्ण करायला हवे असे एक तत्त्व मांडू. संच $X$ चा संचांक $\card{X}$ असा लिहू; मग पुढील अट अपेक्षित आहे:
\[
  \card{A} = \card{B} \text{ iff } \cardeq{A}{B}.
\]
याला आपण \emph{कँटरचे} तत्त्व म्हणू, कारण ही कल्पना कँटर यांच्या मनात बहुधा सर्वप्रथम अगदी स्पष्टपणे होती. (तिचा \emph{ह्यूमच्या} तत्त्वाशी संबंध काय आहे, याविषयी \olref[hp]{sec} मध्ये आणखी सांगू.) म्हणून प्रत्येक $X$ साठी $\card{X}$ ची अशी व्याख्या करणे हे आपले उद्दिष्ट आहे की कँटरचे तत्त्व सिद्ध होईल.

\end{document}
